\chapter{What actually fails}

A poorly run multisig is theater. The contract did its
job. The people did not.

Most protocol losses that involve a ``multisig'' are
not broken signature math. Someone signed the wrong
thing. Someone was unavailable when pause was needed.
Someone put upgrade, treasury, and pause on the same
Safe because it was simpler. Someone copied an address
from an explorer. Someone trusted the UI.

You are the control. The hardware wallet is a locked
box. It will happily sign a theft if you confirm the
hash without reading the effect.

\section{The failure modes that matter}

\begin{itemize}
  \item \textbf{God-mode set.} One Safe owns
  everything. Compromise once, lose the protocol.
  \item \textbf{Rubber stamps.} Threshold signatures
  collected, nobody decoded the calldata.
  \item \textbf{Silent delay.} Timelock queued a
  malicious payload. Nobody was watching, so it
  executed on schedule.
  \item \textbf{Coupled signers.} Same seed, same
  laptop, same apartment, same employer. The M in
  M-of-N was a story.
  \item \textbf{Stuck set.} N-of-N, or keys lost,
  or no backup UI when Safe.app is dark.
  \item \textbf{Lookalike Safe.} Anyone can add you
  as owner of a Safe they created. Dashboard lists
  lie.
\end{itemize}

\section{What this book is not}

It is not a Safe tutorial. Buttons move. The checks
do not: full address, decoded function, simulated
effect, independent hash.

It is not Wallet Security. Seed hygiene, hardware
setup, and personal 2-of-3 live there.

It is not the travel book or the physical-security
book. One sentence each, then go read those:

Do not travel with seed material. A threshold must
still hold if one signer is coerced.

\section{The job}

If you only remember one line from this chapter:

You are not confirming a hash. You are authorizing
an on-chain effect.
